\section{Determinant jets, their short windows, and the actual integral row tests}
\label{sec:determinant-jet-v057}

\status{This chapter proves the full original
\texttt{FormalSeriesBounds.ThmP}, the literal row identity
\texttt{BiValue}, and \texttt{BiContent} and \texttt{BiWindow}.
It then proves the original conversion from a supplied Cartier normal
form to its actual perturbation, including all four fields of
\texttt{MtoHM}, and constructs the integral germs and the exact
change-of-coordinates identity in \texttt{block\_of\_normal\_form}.
The final window uses the same normal form, coordinates and original
row gauge throughout.  Existence of the original normal form
\texttt{RowData.LemmaM} is retained at the last application; it is not
replaced by a new assumption about an unrelated power series.
No whole global obligation is asserted here.}

We want an integral row test whose value is not too divisible by a
prime.  There are two separate issues.  The first is an estimate for
an explicitly perturbed exponential in matrix variables.  The second
is to show that this exponential is the restriction of the prescribed
row gauge, and that its derivatives are values on actual integral
germs.  A scalar estimate by itself would not settle the second issue.

The short-window argument in the first part was prompted by the
adjoint-functor reading edition, revision
\texttt{94107d5f1a355a4ae0bd521123dba9ffd62e00f5},
\texttt{tex/arithmetic.tex}, lines 6432--6528.  We give the coefficient,
factorial and interpolation details.  The remaining parts reconstruct
the original row and coordinate arguments.  All formal-source matches
use revision
\texttt{7cbf7a0c32fca1badf05931879a37c3f8a2b03b4}.
The classical ordinary Cayley identity is the one published input in
the determinant calculation; its exact statement is specified below.

\subsection{Conventions and the local window}

Let $K$ be a characteristic-zero field with an additive
nonarchimedean valuation
$v:K\longrightarrow\mathbb R\cup\{+\infty\}$.
Thus $v(0)=+\infty$, $v(xy)=v(x)+v(y)$, and
$v(x+y)\ge\min(v(x),v(y))$.  We impose neither completeness nor
discreteness.  Fix a prime $p\ge5$ with $v(p)=1$, an odd integer
$h\ge1$, and integers $n\ge1$, $0\le r\le n$.  Put
\begin{gather}
 q=p^h,\qquad \rho=\frac1{p^2-1},\qquad
 \gamma=\frac p{p^2-1},\qquad
 e=1+\frac{q-p}{p^2-1},\qquad c=-p^{-e},\label{djw:parameters}\\
 a=\rho+\frac1{2q}-\frac1{8q},\qquad
 b=\rho+\frac12-\frac1{8q}.
 \label{djw:weights}
\end{gather}
Writing $h=2u+1$ gives
$e=1+p\sum_{j<u}p^{2j}$, so $e$ is a positive integer.

For independent $r$ by $r$ matrices $X,Y$, write $|A|,|B|$ for
the total degrees of multiindices on their entries, and $A!,B!$ for
the products of their factorials.  Let
$U=\sum u_{A,B}X^AY^B$ satisfy
\begin{align}
 v(u_{A,B})+a|A|+b|B|&\ge\frac1{p-1},
       &&\text{for every }A,B,\label{djw:hm1}\\
 u_{A,B}\ne0&\Longrightarrow |B|\ge1.
       \label{djw:hm2}
\end{align}
The second condition says $U(X,0)=0$.  In particular its formal
exponential exists coefficient by coefficient.  All derivatives below
are ordinary derivatives, so a monomial of exponent $A$ contributes
the factorial $A!$, rather than a divided-power normalization.
Define
\begin{equation}
 \Lambda_y=
 \left.\det(\partial_X)^{qy}\det(\partial_Y)^y
 \exp\left(c\sum_{i,j}X_{ij}^{q}Y_{ij}+U(X,Y)\right)
 \right|_{X=Y=0}.
 \label{djw:lambda}
\end{equation}
For a polynomial $P$ and a series $F$ in the same variables, write
\[
 \langle P,F\rangle=P(\partial)F(0)
       =\sum_A P_A A!F_A.
\]
The sum is finite.  Empty determinants and empty products equal one.

% BEGIN DJW1
\paragraph{DJW1: integer and factorial valuations.}
\label{claim:DJW1}
The restriction of $v$ to nonzero integers is the ordinary
$p$-adic valuation.  If a multiindex $A$ has at most $r^2$ coordinates,
$|A|\le M$ and $M\ge1$, then
\begin{equation}
 v(A!)\ge\frac{|A|}{p-1}-r^2(\log_p M+1).
 \label{djw:multifactorial}
\end{equation}
For
\[
 C_r(m)=\prod_{t=1}^{m}\prod_{j=0}^{r-1}(t+j)
       =\prod_{j=0}^{r-1}\frac{(m+j)!}{j!},\qquad C_r(0)=1,
\]
and $r\ge1$, one has
\begin{equation}
 v(C_r(m))\le\frac{rm}{p-1}+r(\log_p r+1).
 \label{djw:capacity}
\end{equation}
Also $v(m!)\le m/(p-1)$ for every $m\ge0$.

\begin{proof}
Repeated addition gives $v(z)\ge0$ for every integer $z$.
If $p\nmid z$, Bezout supplies integers $u,w$ with $uz+wp=1$.
If $v(z)>0$, both terms on the left have positive valuation, a
contradiction.  Removing the power of $p$ from any nonzero integer
proves the first assertion, including its sign.

Counting the multiples of $p,p^2,\ldots$ in a factorial gives
\[
 v(m!)=\sum_{j\ge1}\left\lfloor\frac m{p^j}\right\rfloor
       =\frac{m-s_p(m)}{p-1},
\]
where $s_p(m)$ is its base-$p$ digit sum.  For $m\ge1$,
$s_p(m)\le(p-1)(\log_p m+1)$; $s_p(0)=0$.
Sum this estimate over the coordinates of $A$ to obtain
\eqref{djw:multifactorial}, allowing zero coordinates and $A=0$.
The same formula gives
\[
 v(C_r(m))=
 \frac{rm+\sum_{j<r}s_p(j)-\sum_{j<r}s_p(m+j)}{p-1}.
\]
Discard the last, nonpositive, sum and bound each remaining digit
sum by $(p-1)(\log_p r+1)$.  The upper bound on $v(m!)$ follows
by discarding its nonnegative digit sum.
\end{proof}
% END DJW1

% BEGIN DJW2
\paragraph{DJW2: the determinant factorial and its adjoint identity.}
\label{claim:DJW2}
Let $D=\det X$ with $r\ge1$.  The ordinary Cayley identity is
\begin{equation}
 \det(\partial_X)D^m=
       \prod_{j=0}^{r-1}(m+j)D^{m-1}\qquad(m\ge1).
 \label{djw:cayley}
\end{equation}
Its iteration gives
$\langle D^m,D^m\rangle=C_r(m)$.
If $0\le s\le y$ and $Q$ is homogeneous of degree $rqs$, then
\begin{equation}
 \frac{\langle D^{qy},D^{q(y-s)}Q\rangle}{C_r(qy)}
       =\frac{\langle D^{qs},Q\rangle}{C_r(qs)}.
 \label{djw:adjoint}
\end{equation}

\begin{proof}
We use \eqref{djw:cayley} as the classical polynomial identity in
Caracciolo--Sokal--Sportiello, Theorem 2.1.  It is an identity over
$\mathbb Q$ and hence over every characteristic-zero $K$.
No measure, metric or analytic continuation enters it.

For a single variable, direct comparison of factorials gives
$\langle P,X_jQ\rangle=\langle\partial_jP,Q\rangle$.
Linearity and multiplication give the corresponding identity for
$D$.  Apply it $q(y-s)$ times, and use \eqref{djw:cayley} on the
left polynomial.  Its scalar is $C_r(qy)/C_r(qs)$, which is nonzero
in $K$.  This proves \eqref{djw:adjoint} and, on taking $s=0$,
the stated determinant factorial.
\end{proof}
% END DJW2

% BEGIN DJW3
\paragraph{DJW3: the weighted exponential and the Frobenius error.}
\label{claim:DJW3}
Put $E=\exp U$ and introduce an independent matrix $Z$.
Then
\begin{equation}
 v(E_{A,B})+a|A|+b|B|\ge0,\qquad E(X,0)=1.
 \label{djw:weightedexp}
\end{equation}
For
$\Psi(X,Z)=\det(X^{[q]}-p^eZ)-D^q$, where
$X^{[q]}=(X_{ij}^q)$, every coefficient of a monomial of
$Z$-degree $k$ in $\Psi^s$ is divisible by
\begin{equation}
 p^{ek+(s-k)_+},
 \label{djw:psidiv}
\end{equation}
and its $X$-degree is $q(rs-k)$.

\begin{proof}
A contribution from $U^j/j!$ has weighted valuation at least
$j/(p-1)-v(j!)\ge0$.  At a fixed bidegree only finitely many $j$
occur, because every monomial of $U$ has positive right degree.
The ultrametric inequality proves \eqref{djw:weightedexp}.

Every $Z$ in a determinant expansion supplies a factor $p^e$.
The part of $\Psi$ free of $Z$ is divisible by $p$, since in
characteristic $p$ the Frobenius map commutes with determinants:
$\det(X^{[q]})=(\det X)^q$.
Among $s$ factors of $\Psi$, at most $k$ can have positive
$Z$-degree if the total degree is $k$.  At least $(s-k)_+$ of
them therefore supply an additional $p$.  This proves
\eqref{djw:psidiv} before and after equal monomials are combined.
Every factor has weighted degree $rq$ if an $X$ has weight one
and a $Z$ has weight $q$; the asserted $X$-degree follows.
\end{proof}
% END DJW3

% BEGIN DJW4
\paragraph{DJW4: an exact Newton expansion of the normalized jets.}
\label{claim:DJW4}
For $s\ge0$, let $Q_s(X)$ be the homogeneous part of degree $rqs$
of
$[\Psi(X,\partial_Y)^sE(X,Y)]_{Y=0}$.
Then
\begin{equation}
 \frac{\Lambda_y}{c^{ry}C_r(qy)}
   =\sum_{s=0}^{y}\binom ys a_s,\qquad
 a_s=\frac{\langle D^{qs},Q_s\rangle}{C_r(qs)},\qquad a_0=1.
 \label{djw:newton}
\end{equation}

\begin{proof}
Conjugating $\partial_{Y_{ij}}$ by
$\exp(c\sum X_{ij}^qY_{ij})$ replaces it by
$\partial_{Y_{ij}}+cX_{ij}^q$.
The entries of $X$ commute with these derivatives, and
$c^{-1}=-p^e$.  Consequently the full right operator becomes
\[
 c^{ry}\bigl(D^q+\Psi(X,\partial_Y)\bigr)^y.
\]
Its two summands commute.  Expand by the binomial theorem and set
$Y=0$.  The left polynomial $D^{qy}$ sees only total $X$-degree
$rqy$, so in the term with index $s$ it sees precisely $Q_s$.
Apply DJW2 to remove $D^{q(y-s)}$ from that term.  This is
\eqref{djw:newton}.  Since $E(X,0)=1$, we have $Q_0=1$ and
$a_0=1$.  Every sum used for a fixed coefficient is finite; no
convergence in the valued field is needed.
\end{proof}
% END DJW4

% BEGIN DJW5
\paragraph{DJW5: an explicit lower bound for every late Newton coefficient.}
\label{claim:DJW5}
For $s\ge1$ the coefficients in DJW4 satisfy
\begin{equation}
 v(a_s)\ge
 \frac s8-r^2\bigl(3\log_p r+2\log_p s+h+3\bigr).
 \label{djw:newtonbound}
\end{equation}

\begin{proof}
Fix a term of $Z$-degree $k$ in $\Psi^s$.
Its $X$-degree is $q(rs-k)$, so its contribution to $Q_s$ takes
a coefficient of $E$ of bidegree $(qk,k)$.
Differentiation in $Y$ contributes a multiindex factorial of
degree $k$.  By DJW1 and DJW3 the resulting coefficient has
valuation at least
\[
 ek+(s-k)_++\frac{k}{p-1}
       -r^2(\log_p(rs)+1)-aqk-bk.
\]
The exact identity
\[
 qa+b=e+\frac1{p-1}-\frac{1+q^{-1}}8
\]
reduces this to
\[
 (s-k)_++\frac{k(1+q^{-1})}{8}
       -r^2(\log_p(rs)+1)
 \ge \frac s8-r^2(\log_p(rs)+1).
\]
For $k\le s$ use $s-k+k/8\ge s/8$; for $k\ge s$ use
$k/8\ge s/8$.  Taking finite sums proves this lower bound for
each coefficient of $Q_s$.

The coefficients of $D^{qs}$ are integers.  Its factorial pairing
with $Q_s$ involves only factorials of total degree $rqs$.
DJW1 gives
\[
 v(\langle D^{qs},Q_s\rangle)
 \ge \frac s8-r^2(\log_p(rs)+1)
       +\frac{rqs}{p-1}-r^2(\log_p(rqs)+1).
\]
Subtract the upper bound \eqref{djw:capacity} with $m=qs$.
Since $\log_pq=h$ and $r\le r^2$, the result is
\eqref{djw:newtonbound}.  A zero coefficient or pairing has
valuation $+\infty$ and causes no exception.
\end{proof}
% END DJW5

% BEGIN DJW6
\paragraph{DJW6: the stated numerical threshold is sufficient.}
\label{claim:DJW6}
Let
\begin{equation}
 S(p,h,r)=\left\lceil64r^2
          \bigl(h+3+3\log_p(r+1)\bigr)\right\rceil
          \qquad(r\ge1).
 \label{djw:windowlength}
\end{equation}
The right side of \eqref{djw:newtonbound} is strictly positive
for every integer $s\ge S(p,h,r)$.

\begin{proof}
Set $A=h+3+3\log_p(r+1)\ge4$ and $s_0=64r^2A$.
The function
\[
 F(t)=\frac{t}{8r^2}-h-3-3\log_p r-2\log_p t
\]
has positive derivative for $t\ge s_0$:
$1/(8r^2)-2/(t\log p)>0$ there.
With $b_0=\log_p(r+1)$, direct substitution gives
\begin{align*}
 F(s_0)&\ge7A-4b_0-2\log_p64-2\log_pA\\
       &>7A-\frac{4A}{3}-6-2A
        =\frac{11A}{3}-6>0.
\end{align*}
We used $b_0\le A/3$, $\log_p64<3$ for $p\ge5$, and
$\log_pA\le A$ for $A\ge1$.  Monotonicity and the ceiling
now prove the assertion.
\end{proof}
% END DJW6

% BEGIN DJW7
\paragraph{DJW7: integral interpolation on every translated interval.}
\label{claim:DJW7}
Let $K$ be any field with an additive nonarchimedean valuation,
$S\ge1$, and $a_s\in K$ with $a_0=1$ and $v(a_s)>0$ for
every $s\ge S$.  For each $Y_0\ge0$, some
$Y_0\le y<Y_0+S$ satisfies
\[
 v\left(\sum_{s=0}^{y}\binom ys a_s\right)\le0.
\]
This assertion does not require characteristic zero or a normalized
prime.

\begin{proof}
The integer matrix
$V_{ij}=\binom{Y_0+i}{j}$, $0\le i,j<S$, factors as $BC$,
where $B_{ik}=\binom ik$ is lower triangular with diagonal one
and $C_{kj}=\binom{Y_0}{j-k}$ is upper triangular with diagonal
one.  The factorization is the elementary binomial convolution.
Thus $V^{-1}$ has integer entries, over $\mathbb Z$ and after
mapping to $K$.  Every integer has nonnegative valuation.

If all $S$ values in question had positive valuation, their finite
tails with $s\ge S$ would also have positive valuation.  Subtract
those tails, treating $\binom ys$ as zero when $s>y$.
The resulting vector of truncated sums is
$V(a_0,\ldots,a_{S-1})^t$ and has every entry of positive valuation.
Multiplication by $V^{-1}$ preserves this property, which would
give $v(a_0)>0$.  This contradicts $v(1)=0$.
\end{proof}
% END DJW7

% BEGIN DJW8
\paragraph{DJW8: the full original short-window theorem.}
\label{claim:DJW8}
Under \eqref{djw:parameters}--\eqref{djw:hm2}, if $r\ge1$,
each $Y_0\ge0$ admits $Y_0\le y<Y_0+S(p,h,r)$ with
\begin{equation}
 v(\Lambda_y)\le
 ry\bigl((q+1)\gamma-1\bigr)+n(\log_p n+1).
 \label{djw:thmp}
\end{equation}
If $r=0$, then $\Lambda_y=1$ for every $y\ge0$.
This is the complete original
\texttt{FormalSeriesBounds.ThmP}, including its zero-size conjunct.

\begin{proof}
For $r\ge1$, use DJW4--DJW7 to find a normalized value with
valuation at most zero.  DJW1 then gives
\begin{align*}
 v(\Lambda_y)
 &\le-ery+\frac{rqy}{p-1}+r(\log_pr+1)\\
 &=ry\bigl((q+1)\gamma-1\bigr)+r(\log_pr+1)\\
 &\le ry\bigl((q+1)\gamma-1\bigr)+n(\log_pn+1).
\end{align*}
The middle identity is $q/(p-1)-e=(q+1)p/(p^2-1)-1$.
The last inequality follows from $1\le r\le n$ and monotonicity
of the two nonnegative factors.

For $r=0$ there are no variables.  The only monomial has right
degree zero, so \eqref{djw:hm2} forces $U=0$.  Both determinant
operators are identities and the exponential is one.  In particular
we do not try to use a zero-length interval in this case.
\end{proof}
% END DJW8

\subsection{The exact row identity}

The next algebraic statements have a wider scope than DJW8.
In this subsection $K$ is any characteristic-zero field, and the
integers $l,i,n,r,y$ are arbitrary nonnegative integers with $r\le n$.
There is no prime, valuation, oddness or separated-row hypothesis.
For these general indices use the original integer
\[
 e_{l,i}=1+l\sum_{j<(i-1)/2}(l^2)^j,
 \qquad q=l^i,\qquad c=-\bigl(l^{e_{l,i}}\bigr)^{-1},
\]
where the sum uses integer truncation and the field inverse is zero
at zero.  This agrees with \eqref{djw:parameters} in the prime,
odd-distance case.  Let $Z,W$ be independent $n$ by $n$ matrices.
Restriction $\operatorname{res}_r$ sets every variable outside the
leading $r$ by $r$ blocks to zero.

For completeness the source defines a truncated coefficient operation
on every series $H$, even one with nonzero constant term, by
\[
 \operatorname{E}(H)_A
       =\sum_{j=0}^{|A|}[X^A]H^j/j!.
\]
When $H(0)=0$ this is the ordinary formal exponential.  We use
the symbol $\operatorname{E}$ only to state its unrestricted
restriction law precisely.

% BEGIN DJW9
\paragraph{DJW9: restriction, support and genuine formal exponentials.}
\label{claim:DJW9}
Restriction is a ring homomorphism and
\[
 \operatorname{res}_r\operatorname{E}(H)
       =\operatorname{E}(\operatorname{res}_rH)
\]
for every $H$.  If $H(0)=H'(0)=0$, then
$\exp(H+H')=\exp(H)\exp(H')$.
Series supported on a specified subset of variables remain so under
sum, product, powers and $\operatorname{E}$.  A substitution by
series with zero constant term preserves order of vanishing and
commutes with ordinary formal exponentiation of a zero-constant
series.

\begin{proof}
The coefficient of a product is the finite sum over decompositions
of its exponent.  An exponent supported in the retained variables
can be decomposed only into exponents supported there.  This proves
multiplicativity.  The retained exponent has the same degree in
the small and large variable sets, so the upper limit $|A|$ in
$\operatorname{E}$ is unchanged.  Applying multiplicativity to
each $H^j$ proves the unrestricted restriction law.

For zero-constant series, the coefficient of degree at most $N$
in a power with exponent $>N$ is zero.  The usual binomial identity,
divided by its factorial, therefore proves the exponential product
identity by finite sums in every degree.  The assertions about
support follow from the same coefficient decompositions.

Under a zero-constant substitution a monomial of degree $d$ has
no terms of degree less than $d$.  Hence to compute a coefficient
of degree $N$ one may first truncate the substituted series and
the exponential at degree $N$.  Both are then finite polynomial
identities.  This proves the last two assertions without any
topological hypothesis on $K$.
\end{proof}
% END DJW9

% BEGIN DJW10
\paragraph{DJW10: determinant coefficients with prescribed diagonal rows.}
\label{claim:DJW10}
Every monomial of $(\det Z)^A(\det W)^B$ has degree $A$ in
each row of $Z$ and $B$ in each row of $W$.
If in every row $\rho\ge r$ a monomial uses only the diagonal
entries, with exponents $A$ and $B$, then its coefficient is the
coefficient of its remaining exponent in
$(\det Z_{<r,<r})^A(\det W_{<r,<r})^B$.
In particular this coefficient vanishes if the remaining monomial
uses a column $\ge r$.

\begin{proof}
Each permutation monomial of a determinant uses exactly one entry
in each row and column.  Taking products proves the degree assertion.
Set all off-diagonal entries in rows $\rho\ge r$ equal to zero.
The resulting matrix is block upper triangular, with its lower
diagonal entries on the diagonal of the lower block.  Its determinant
is the leading-block determinant times their product, with sign $+1$.
Raise to the stated powers and compare coefficients.
Specialization kills precisely the monomials using a deleted variable,
so it does not change any coefficient under consideration.  This
argument proves the source's full \texttt{DetPowDiag} over every
commutative coefficient ring, including $r=0$, $r=n$ and zero powers.
\end{proof}
% END DJW10

% BEGIN DJW11
\paragraph{DJW11: the full diagonal-block jet formula.}
\label{claim:DJW11}
Suppose $G$ uses only rows $\rho<r$ and $H$ only rows $\rho\ge r$.
Assume that, among exponents having degree at most $A$ in each
right-hand $Z$ row and at most $B$ in each right-hand $W$ row,
the only nonzero coefficient of $H$ is
\[
 \prod_{\rho=r}^{n-1}Z_{\rho\rho}^{A}W_{\rho\rho}^{B},
 \qquad\text{with coefficient }1.
\]
Then
\begin{equation}
 \langle(\det Z)^A(\det W)^B,GH\rangle
  =(A!B!)^{n-r}
    \langle(\det Z_{<r,<r})^A(\det W_{<r,<r})^B,
                \operatorname{res}_rG\rangle.
 \label{djw:blockjet}
\end{equation}
In particular this proves the full source \texttt{BlockJet} on taking
$A=Ly$, $B=y$ for arbitrary $L,y\ge0$.

\begin{proof}
For a monomial in the differential polynomial, DJW10 fixes each
row degree.  In the coefficient convolution of $GH$, the exponent
on the last rows comes entirely from $H$.  It is therefore within
the stated row-degree limits.  The hypothesis on $H$ forces exactly
the displayed diagonal exponent, with coefficient one.
DJW10 then identifies the polynomial coefficient on the first rows
with its leading-block coefficient and discards any remaining extra
columns.  The factorial of the complete exponent is the factorial
of its first-block exponent times $(A!B!)^{n-r}$, since the supports
are disjoint.  Summing gives \eqref{djw:blockjet}.
This coefficient proof also covers empty blocks and $A=0$ or $B=0$;
no differentiation by a positive number of variables was presumed.
\end{proof}
% END DJW11

Let $H_\rho,T_\rho,T'_\rho\in K[[Z,W]]$ satisfy the five
original \texttt{BiHyp} requirements:
each $H_\rho$ uses only row $\rho$ and has zero constant term;
for $\rho\ge r$, $T_\rho$ uses only that row's $Z$ variables
and equals $Z_{\rho\rho}$ to degree one, while $T'_\rho$
uses only that row's $W$ variables and equals $W_{\rho\rho}$
to degree one; finally
\begin{equation}
 \operatorname{res}_r\sum_{\rho<r}H_\rho
       =c\sum_{a,b<r}Z_{ab}^{q}W_{ab}+U.
 \label{djw:biR3}
\end{equation}
Here ``equals to degree one'' includes equality of the constant
and every linear coefficient, not just the indicated diagonal
coefficient.  Define the literal row product
\[
 F_y=\prod_{\rho=0}^{n-1}
       \left(\begin{cases}1,&\rho<r,\\
           T_\rho^{qy}(T'_\rho)^y,&\rho\ge r
       \end{cases}\right)\exp(H_\rho).
\]

% BEGIN DJW12
\paragraph{DJW12: the literal row-value identity \texttt{BiValue}.}
\label{claim:DJW12}
For every $y\ge0$ and all the algebraic indices just specified,
\begin{equation}
 \langle(\det Z)^{qy}(\det W)^y,F_y\rangle
       =((qy)!y!)^{n-r}\Lambda_y.
 \label{djw:bivalue}
\end{equation}
The $\Lambda_y$ on the right is the original small-block series
formed from the same $U$ in \eqref{djw:biR3}; neither of the
valuation hypotheses \eqref{djw:hm1}--\eqref{djw:hm2} is needed
for this identity.

\begin{proof}
Factor $F_y=GH$, where
$G=\exp(\sum_{\rho<r}H_\rho)$ and $H$ is the product over
the remaining rows.  DJW9 justifies this factorization and its
row supports.  A nonzero monomial in $T_\rho^{qy}$ has
$Z$-row degree at least $qy$, with equality only for
$Z_{\rho\rho}^{qy}$ and coefficient one.  This follows by choosing
the linear term from every factor; every other term has larger degree.
The same statement holds for $(T'_\rho)^y$.
Every monomial of $\exp(H_\rho)$ has nonnegative degrees on both
sides, and its constant coefficient is one.  Thus, within the degree
limits $qy,y$, the product on the last rows has exactly the diagonal
coefficient required by DJW11.  A nonconstant exponential term cannot
survive these limits.  This remains valid if either power is zero.

Apply DJW11 and then DJW9.  The restricted first-row factor is
\[
 \operatorname{res}_rG
   =\exp\left(\operatorname{res}_r\sum_{\rho<r}H_\rho\right)
   =\exp\left(c\sum_{a,b<r}Z_{ab}^{q}W_{ab}+U\right).
\]
Its exponent has zero constant term because the left side of
\eqref{djw:biR3} does.  The remaining jet is precisely $\Lambda_y$.
This proves every quantifier of \texttt{BiValue}, including $r=0$,
$n=0$, $l=0$ and $i=0$ when their algebraic data are supplied.
\end{proof}
% END DJW12

Return now to the valued prime and odd-distance hypotheses of the
first subsection.  For brevity write
\[
 B_r(y)=ry((q+1)\gamma-1)+n(\log_p n+1).
\]

% BEGIN DJW13
\paragraph{DJW13: the exact factorial content loss.}
\label{claim:DJW13}
If $v(\Lambda)\le B_r(y)$, then
\begin{equation}
 v\bigl(((qy)!y!)^{n-r}\Lambda\bigr)
       \le B_n(y)+\frac{(n-r)(q+1)y}{p-1}.
 \label{djw:content}
\end{equation}
This is the full original \texttt{BiContent} for every scalar
$\Lambda$ satisfying its bound.

\begin{proof}
DJW1 bounds the two factorial valuations by $qy/(p-1)$ and
$y/(p-1)$.  Also
\[
 (q+1)\gamma-1\ge0,
\]
because $q\ge p$ and
$(p+1)p/(p^2-1)>1$.  Hence $B_r(y)\le B_n(y)$.
Add these inequalities using multiplicativity of the valuation.
The assumed finite upper bound also excludes $\Lambda=0$;
there is no subtraction of an infinite valuation.
\end{proof}
% END DJW13

% BEGIN DJW14
\paragraph{DJW14: \texttt{BiWindow} and \texttt{content\_and\_window}.}
\label{claim:DJW14}
Suppose $1\le r\le n$, $n-r\le(h+1)m$, and the literal
\texttt{BiHyp} data satisfy \eqref{djw:hm1}--\eqref{djw:hm2}
for their same $U$.  Every $Y_0\ge0$ has an index
$Y_0\le y<Y_0+S(p,h,r)$ such that
\begin{equation}
 v\bigl(\langle(\det Z)^{qy}(\det W)^y,F_y\rangle\bigr)
 \le B_n(y)+\frac{(h+1)m(q+1)y}{p-1}.
 \label{djw:biwindow}
\end{equation}
Together with DJW13, this supplies the full original
\texttt{content\_and\_window} theorem.

\begin{proof}
Choose $y$ by DJW8 for the small block.  Substitute the exact
identity DJW12 and apply DJW13.  The final comparison replaces
$n-r$ by $(h+1)m$ in a nonnegative expression.  This is exactly
the source's \texttt{biWindow\_of} implication with both its
inputs now proved; it adds no geometric premise.
\end{proof}
% END DJW14

\subsection{The perturbation attached to the original normal form}

We now fix the original datum
$\mathcal X=\texttt{RowData}(K,V,d,I,J)$.
Its field is algebraically closed of characteristic zero, $d>0$,
and $V=(v,p)$ is the original valued-prime datum.
There are two finite sets of tangent labels $I_0,I_1$, their
Cartier data $G_0,G_1$, and the original leaf gauges
\[
 H_{g,\alpha,\alpha'}\in
 K[[z_w:w\in I_0\amalg I_1]].
\]
The datum contains the assigned types, the ranks, integral pairings,
Cartier identities, and comparison of integral-germ spaces between
coordinate systems.  None of those fields is silently changed here.
We use in particular its axis-vanishing property: a gauge coefficient
supported on just one of the two sides is zero.  Thus each such
gauge has zero constant term.

The actual integers in this datum are denoted $n,j,j'$ and
$m=j+j'$.  Fix $g\in\{0,1\}$, a type $\alpha\in\mathbb Z/d\mathbb Z$,
and $h\ge0$.  Until DJW21, these algebraic constructions do not
require $h$ odd or $p\ge5$.
Put $q=p^h$ and use the original integer $e_{p,h}$ above.
For row vectors $u_c$ on $I_g$ and $v_b$ on $I_{1-g}$ write
\[
 H(u,v)(x,y)=H_{g,\alpha,\alpha+h}
       \left(\sum_cu_c x_c,\sum_bv_b y_b\right).
\]
This is the literal linear pullback of the given gauge, not a
replacement series chosen for the estimate.

Here is the exact normal-form premise used below.  The given arrays
$U_0,U'_0$ each have $n$ rows.  An original
\texttt{MConcl} consists of $0\le r\le n$, vectors $e_c$ for
$c<r$, vectors $v'_b$ for $b<n$, and a matrix $\Phi$ of size
$n$ by $r$, with the following properties:
\begin{enumerate}
\item $n\le r+(h+1)j'+j$.
\item The $e_c$ are integral and have an integral right inverse:
there is $w_{x,c}$ with
$\sum_x e_{c,x}w_{x,c'}=\delta_{cc'}$.
\item $e=A U_0$ for an integral $r$ by $n$ matrix $A$.
\item $v'=Y U'_0$ for an integral $n$ by $n$ matrix $Y$ with
$v(\det Y)=0$.
\item $\Phi$ is integral and its first $r$ rows are $-I_r$.
\item Every coefficient $R_{A,B}$ of
\begin{equation}
 R(x,y)=H(e,v')(x,y)
        -p^{-e_{p,h}}\sum_{b<n,c<r}\Phi_{bc}x_c^q y_b
 \label{djw:actualremainder}
\end{equation}
satisfies
$v(R_{A,B})+a|A|+b|B|\ge1/(p-1)$, with weights
\eqref{djw:weights} for the present $p,h$.
\item $R_{A,0}=0$ for every $A$.
\end{enumerate}
These are exactly the original fields, with its first pair of
integer inequalities separated for readability.  When $h=1$ the
additional original \texttt{MPlusConcl} means
$v(R_{A,B})\ge0$ whenever $|A|<p^2$ and $|B|<p$.
Saturation here is the displayed integral right inverse; no
Noetherian or discrete-valuation interpretation is substituted.

% BEGIN DJW15
\paragraph{DJW15: all four original perturbation conclusions.}
\label{claim:DJW15}
Restrict \eqref{djw:actualremainder} by setting $y_b=0$ for
$b\ge r$ and denote the result by $R_0(x,y)$.
For independent $r$ by $r$ matrices $Z,W$ put
\begin{equation}
 U(Z,W)=\sum_{\rho<r}R_0(Z_{\rho,*},W_{\rho,*}).
 \label{djw:actualU}
\end{equation}
This is exactly
\texttt{restrictedCurveSystemSeries} of the original normal form.
It satisfies all of \texttt{MtoHM}:
\begin{align}
 \sum_{\rho<r}H(e,v'_{<r})(Z_{\rho,*},W_{\rho,*})
   &=-p^{-e_{p,h}}\sum_{\rho,c<r}Z_{\rho c}^qW_{\rho c}+U,
        \label{djw:actualmain}\\
 v(U_{A,B})+a|A|+b|B|&\ge\frac1{p-1},\qquad U_{A,0}=0.
        \label{djw:actualHM}
\end{align}
If $h=1$ and \texttt{MPlusConcl} holds, the decomposition
\eqref{djw:actualU} has each summand supported in one row and
integral in all bidegrees $|A|<p^2$, $|B|<p$.

\begin{proof}
The restriction is the substitution $y_b\mapsto0$ for $b\ge r$.
It commutes with the actual linear pullback of the gauge by DJW9's
coefficient argument.  Since the first $r$ rows of $\Phi$ are
$-I_r$, it takes the explicit main term in
\eqref{djw:actualremainder} to
$-p^{-e_{p,h}}\sum_{c<r}x_c^qy_c$.
Thus the restricted remainder is precisely the difference defining
the original restricted series.  Substitute each row and sum to
obtain \eqref{djw:actualmain}, including its negative sign.

A coefficient of $R_0$ is exactly a coefficient of $R$ with zeros
inserted in the deleted positions.  Its left and right total degrees
are unchanged.  Hence it retains the weighted inequality and
the first-axis vanishing of $R$.  Row substitution is an injective
renaming of variables.  A coefficient in its image retains those
same two degrees; a coefficient outside its image is zero.
For a fixed matrix exponent, every summand in
\eqref{djw:actualU} therefore has valuation at least the same
number $1/(p-1)-a|A|-b|B|$.  The valuation inequality for finite
sums proves \eqref{djw:actualHM}, even where several summands
contribute or cancel.  First-axis coefficients vanish in every
summand, proving the second assertion without relying on such
cancellation.

When $h=1$, use the additional coefficient bound on $R$ before
restriction.  The two degree cutoffs are again unchanged by
restriction and injective row renaming.  This proves the exact
\texttt{HSeparatedRowSum} conjunct, with the displayed summands
as its witnesses.  If $r=0$, the sum is empty and every assertion
has the same literal meaning.  We have imposed neither the
alternating-type condition nor oddness in proving this conversion.
\end{proof}
% END DJW15

\subsection{Constructing the actual integral row tests}

Write $\mathcal O=\{a\in K:v(a)\ge0\}$.
The integral germs of a coordinate system $t=(t^g_x)$ are, by
the original definition, the series
\begin{equation}
 \gamma\bigl((t^g_x)_{g,x}\bigr),\qquad
       \gamma\in\mathcal O[[z_w:w\in I_0\amalg I_1]],
 \label{djw:intgerm}
\end{equation}
where each $t^g_x$ is inserted into its own side of the disjoint
variable set.  This definition does not assert that the coefficients
of $t^g_x$ in logarithmic coordinates are themselves integral.

% BEGIN DJW16
\paragraph{DJW16: completing split integral rows without a discrete valuation.}
\label{claim:DJW16}
Suppose $A\in M_{r,n}(\mathcal O)$,
$B\in M_{n,r}(\mathcal O)$ and $AB=I_r$.
For every $r\le j\le n$ there are integral matrices
$A_j\in M_{j,n}(\mathcal O)$ and $B_j\in M_{n,j}(\mathcal O)$
with $A_jB_j=I_j$ and first $r$ rows of $A_j$ equal to $A$.
For $j=n$ this gives inverse square matrices
$\bar A,\bar B\in M_n(\mathcal O)$ and
$v(\det\bar A)=v(\det\bar B)=0$.

\begin{proof}
It suffices to add one row to a split pair with $j<n$.
Over $K$, the kernel of $A_j:K^n\to K^j$ contains a nonzero
vector $u$.  Choose an index $k$ at which $v(u_k)$ is minimal.
The finite minimum is finite, so $u_k\ne0$.
Set $s_x=u_x/u_k$.  Every $s_x$ is integral, $s_k=1$, and
$A_js=0$.  Define the new row
\[
 r_x=\delta_{kx}-\sum_{c<j}(B_j)_{kc}(A_j)_{cx}.
\]
It is integral.  Direct multiplication gives
\[
 rB_j=0,\qquad rs=1,\qquad A_js=0.
\]
Append $r$ to $A_j$ and append $s$ as a column to $B_j$.
These three identities together with $A_jB_j=I_j$ prove that
the enlarged matrices still multiply to the identity.
Induction supplies every intermediate size, including the empty
initial set of rows.

For $j=n$, taking determinants gives
$\det\bar A\det\bar B=1$.
Both determinants are integral, so both have valuation zero;
the matrices are inverses over $K$ and already have integral
entries.  The proof used only a minimum in a finite list of
valuations.  No uniformizer, completeness or residue-field
finiteness is needed.  This proves the original
\texttt{complete\_rows}, with all its intermediate sizes.
\end{proof}
% END DJW16

% BEGIN DJW17
\paragraph{DJW17: actual standard Cartier coordinates.}
\label{claim:DJW17}
For either original Cartier datum $G_g$ one can choose its system
$S_g$ and its associated formal coordinates $t^g_x$ so that
\[
 (S_g).\mathrm{IsT}(t^g),\qquad
 t^g_x=z_x+\text{terms of total degree at least two}.
\]
These choices depend on the datum, not on the later integer $y$.

\begin{proof}
The original Cartier axiom supplies a typed curve with any integral
typed initial tangent vector.  Apply it to each standard basis vector.
The resulting system has identity tangent matrix, hence unit
determinant.  If $\alpha_{c,s,x}$ is its tangent-side coefficient
at Cartier index $s$, then $\alpha_{c,0,x}=\delta_{cx}$.
The coordinate equations are the formal equalities
\begin{equation}
 \sum_{s\ge0,c}\alpha_{c,s,x}(t_c)^{p^s}=z_x.
 \label{djw:cartiercoord}
\end{equation}
For any fixed degree the sum is finite, since $p\ge2$ and the
coordinates have zero constant term.  Separate the $s=0$ terms
and write this as $t_x=z_x-Q_x(t)$, where every monomial of
$Q_x$ has degree at least two.

Starting from zero, iterate $t\mapsto z-Q(t)$.
If two zero-constant tuples agree below degree $N$, their images
agree below degree $N+1$.  Indeed each difference
$u^a-v^a$ factors by $u-v$, and its remaining $a-1\ge1$
factors have positive degree.  Thus every coefficient of the
iterates stabilizes after finitely many steps.  The stabilized
tuple solves \eqref{djw:cartiercoord}, and the same argument
proves uniqueness among zero-constant tuples.  Its degree-one
part is $z$.

Taking coefficients in \eqref{djw:cartiercoord} gives precisely
the source's finite equation \texttt{Sys.IsT}; terms with
$p^s$ above the coefficient degree vanish.  This constructs the
actual systems and coordinates rather than merely positing a
formal change of variables.
\end{proof}
% END DJW17

% BEGIN DJW18
\paragraph{DJW18: linear substitution and the ordinary differential pairing.}
\label{claim:DJW18}
Let $L:z\mapsto w$ be a linear substitution between finite sets
of formal variables over a characteristic-zero field.  For every
polynomial $P$ in the $z$ variables and every series $F$ in the
$w$ variables,
\begin{equation}
 \langle P,F\circ L\rangle
       =\langle P(L^t\partial_w),F\rangle.
 \label{djw:linearjet}
\end{equation}
Here the right side means the polynomial obtained by replacing
each differential letter $\partial_{z_j}$ by
$\sum_iL_{ij}\partial_{w_i}$.  If, row by row, a square matrix
of these letters becomes $P_0A^t$, its determinant becomes
$\det A\det P_0$.

\begin{proof}
For one derivative the coefficient formula for substitution gives
the ordinary chain rule
$\partial_{z_j}(F\circ L)
=\sum_iL_{ij}(\partial_{w_i}F)\circ L$.
It is enough to check this on monomials; the factorial coefficients
agree directly.  Iterate the identity and set $z=0$, then extend
linearly over $P$.  In every prescribed degree all calculations
involve finite polynomials, so the argument applies to formal
series.  Determinant multiplicativity gives the last assertion.
In particular the scalar is $\det A$, not its inverse.  This
orientation will fix the exact scalar in the actual row value.
\end{proof}
% END DJW18

For the rest of the construction take injective maps
$\iota:[n]\hookrightarrow I_g$ and
$\iota':[n]\hookrightarrow I_{1-g}$, and specialize the
given normal form to
\[
 (U_0)_{b,x}=\delta_{x,\iota(b)},\qquad
 (U'_0)_{b,x}=\delta_{x,\iota'(b)}.
\]
These are the same selected letters as in the original
\texttt{block\_of\_normal\_form} theorem.  No further
type or alternating hypothesis is needed once this original
normal form is supplied.

% BEGIN DJW19
\paragraph{DJW19: the actual row substitution and all five \texttt{BiHyp} fields.}
\label{claim:DJW19}
For the given normal form, there are systems $S_g$, coordinates
$t^g$, row exponents $H_\rho$ and row coordinates
$T_\rho,T'_\rho$ satisfying the original \texttt{BiHyp} with
the exact perturbation \eqref{djw:actualU}.
All these objects are chosen before $y$.

\begin{proof}
Let $w$ be the integral right inverse of $e$ in the normal form.
The matrix $B_{b,c}=w_{\iota(b),c}$ satisfies $AB=I_r$:
substitute $e_{c,x}=\sum_bA_{cb}\delta_{x,\iota(b)}$ in
$\sum_xe_{c,x}w_{x,c'}=\delta_{cc'}$.
Apply DJW16 to obtain square inverse matrices $\bar A,\bar B$
extending the first $r$ rows of $A$.
Also put $Z_0=Y^{-1}$.  It is integral because the adjugate of
the integral matrix $Y$ is integral and $v(\det Y)=0$.
Define
\[
 E_{c,x}=\sum_b\bar A_{cb}\delta_{x,\iota(b)}.
\]
Its first $r$ rows are $e$.  Moreover
$E_{c,\iota(b)}=\bar A_{cb}$ and
$v'_{c,\iota'(b)}=Y_{cb}$, by injectivity of the selected letters.
Choose the standard systems and coordinates from DJW17.

For each $\rho<n$, the actual row substitution is
\begin{align*}
 z_{(g,x)}&\longmapsto\sum_cE_{c,x}Z_{\rho c},\\
 z_{(1-g,x)}&\longmapsto\sum_cv'_{c,x}W_{\rho c}.
\end{align*}
Call it $L_\rho$, and define
\begin{align*}
 H_\rho&=H_{g,\alpha,\alpha+h}\circ L_\rho,\\
 T_\rho&=
 \left(\sum_b\bar B_{b\rho}t^g_{\iota(b)}\right)
       \left(z_x=\sum_cE_{c,x}Z_{\rho c}\right),\\
 T'_\rho&=
 \left(\sum_b(Z_0)_{b\rho}t^{1-g}_{\iota'(b)}\right)
       \left(z_x=\sum_cv'_{c,x}W_{\rho c}\right).
\end{align*}
Every substituted variable has zero constant term, so these formal
substitutions are defined.

The first series uses only row $\rho$ and has constant term zero,
by the axis-vanishing of the original gauge.  The second uses only
the left side of that row.  Its degree-one part is
\[
 \sum_{b,c}\bar B_{b\rho}E_{c,\iota(b)}Z_{\rho c}
       =\sum_c(\bar A\bar B)_{c\rho}Z_{\rho c}
       =Z_{\rho\rho}.
\]
Every higher term of $t$ remains of degree at least two under
linear substitution.  The analogous computation with $YZ_0=I$
gives the precise right diagonal coordinate.

Finally restrict $\sum_{\rho<r}H_\rho$ to the leading blocks.
The first $r$ rows of $E$ are $e$, and the surviving right
columns are exactly $v'_{<r}$.  Thus the restricted sum is the
left side of \eqref{djw:actualmain}.  DJW15 identifies it with
the required main term plus the very same $U$.
This proves all five \texttt{BiHyp} fields, including the exact
restriction equality.  No choice in this construction depends
on a derivative order or a window origin.
\end{proof}
% END DJW19

To state the final value identity, use independent ambient variables
$z_{\rho,w}$ for $\rho<n$ and $w\in I_0\amalg I_1$.
Let
\[
 D_g=\det(z_{\rho,(g,\iota(c))})_{\rho,c<n},\qquad
 D_{1-g}=\det(z_{\rho,(1-g,\iota'(c))})_{\rho,c<n},
 \qquad P_y=D_g^{qy}D_{1-g}^{y}.
\]
These are the original selected-letter determinants.  The actual
row value for germs $f_\rho$ is
\begin{equation}
 \mathcal V_y(f)=
 \left\langle P_y,
  \prod_{\rho<n}f_\rho(z_{\rho,*})
              \exp H_{g,\alpha,\alpha+h}(z_{\rho,*})\right\rangle.
 \label{djw:actualvalue}
\end{equation}
Thus its row-kind is the same
$\operatorname{some}(g,\alpha,\alpha+h)$ in every row.

% BEGIN DJW20
\paragraph{DJW20: integral germs and the exact original value comparison.}
\label{claim:DJW20}
For every $y\ge0$, there are actual germs
$f_{\rho,y}\in\operatorname{IntGerm}(V,t)$ such that
\begin{equation}
 \langle(\det Z)^{qy}(\det W)^y,F_y\rangle
   = (\det\bar A)^{qy}(\det Y)^y\,
                     \mathcal V_y(f_{*,y}).
 \label{djw:exactscalar}
\end{equation}
The scalar has valuation zero.  In particular
\begin{equation}
 v(\mathcal V_y(f_{*,y}))
 \le v\bigl(\langle(\det Z)^{qy}(\det W)^y,F_y\rangle\bigr),
 \label{djw:actualcomparison}
\end{equation}
with equality for the displayed witnesses.  DJW19 and this statement
are the full original \texttt{block\_of\_normal\_form} conclusion
for its supplied normal form and restriction identity.
More generally, for every supplied $h_1$ whose main-term expression
equals the same restricted gauge, the construction satisfies
\texttt{BiHyp} and \eqref{djw:exactscalar} with that $h_1$;
no weighted estimate on $h_1$ is needed for this statement.

\begin{proof}
Define integral polynomials on the original tangent-label variables by
\[
 \gamma_{\rho,y}(z)=
 \begin{cases}
 1,&\rho<r,\\[2pt]
 \displaystyle
 \left(\sum_b\bar B_{b\rho}z_{(g,\iota(b))}\right)^{qy}
 \left(\sum_b(Z_0)_{b\rho}z_{(1-g,\iota'(b))}\right)^y,
      &\rho\ge r.
 \end{cases}
\]
All their coefficients lie in $\mathcal O$: their linear
coefficients do, and this ring is closed under finite sums and
products.  Put $f_{\rho,y}=\gamma_{\rho,y}(t)$ with the same
coordinate system chosen in DJW19.  Formula \eqref{djw:intgerm}
proves their required integrality in the original sense.

Apply the row substitutions $L_\rho$ simultaneously to the series
inside \eqref{djw:actualvalue}.  By construction each
$f_{\rho,y}$ becomes the corresponding coordinate powers in $F_y$,
and its gauge becomes $H_\rho$.  DJW9 permits composition of the
zero-constant substitutions and commutation with the exponential.
Thus the substituted product is exactly $F_y$.

Apply DJW18 to its derivative polynomial.
The matrix of left derivative letters is the selected-letter
matrix multiplied on the right by $\bar A^t$, and the right
one is multiplied on the right by $Y^t$.  Their determinants
therefore contribute $\det\bar A$ and $\det Y$, respectively.
Taking the powers $qy$ and $y$ proves \eqref{djw:exactscalar}.
Both determinants have valuation zero, by DJW16 and the normal-form
premise.  Taking valuations gives the equality and hence
\eqref{djw:actualcomparison}, even when the value is zero.

For the literal source theorem, the supplied series $h_1$ need
only satisfy its displayed restriction identity, rather than be
named by \eqref{djw:actualU}.  The construction just given uses
that identity in DJW19 and no other property of $h_1$.
Consequently it proves the source's full statement for every such
$h_1$, not only for our canonical choice.  The canonical choice
is used in DJW21 to obtain the weighted estimates from DJW15.
\end{proof}
% END DJW20

% BEGIN DJW21
\paragraph{DJW21: a short window on the actual theta-row data.}
\label{claim:DJW21}
Assume now $p\ge5$, $h$ odd and $n\ge1$, and keep the original
normal form and selected letters of DJW19.
There is one choice of the original systems $S$ and coordinates
$t$, before any $y$ or $Y_0$, with the following property.
Zero is an admissible order, and every $Y_0\ge0$ has
$Y_0\le y<Y_0+S(p,h,n)$ and integral germs $f_{\rho,y}$ with
\begin{equation}
 v(\mathcal V_y(f_{*,y}))\le
 \begin{cases}
 0,&y=0,\\[2pt]
 \displaystyle
 ny((q+1)\gamma-1)+n(\log_p n+1)
       +\frac{(h+1)m(q+1)y}{p-1},&y>0.
 \end{cases}
 \label{djw:actualwindow}
\end{equation}
This is the original \texttt{Windows} conclusion on the actual
selected-letter polynomial and its exact \texttt{blockBd},
conditional only on the stated original normal-form premise.

\begin{proof}
Use DJW15 to supply both weighted hypotheses for the exact $U$,
and DJW19--DJW20 for the actual coordinate and germ comparison.
The normal-form rank inequality gives
\[
 n-r\le(h+1)j'+j\le(h+1)(j+j')=(h+1)m.
\]
If $r\ge1$, DJW14 supplies a suitable $y$ in a window of length
$S(p,h,r)$.  This is at most $S(p,h,n)$, since $r^2$ and
$h+3+3\log_p(r+1)$ are nonnegative increasing functions of $r$.
Combine its estimate with \eqref{djw:actualcomparison}.

If $r=0$, take $y=Y_0$ itself.  DJW8 gives $\Lambda_y=1$,
and DJW12--DJW13 give the same bound after the factorial loss.
The required interval has positive length because $n\ge1$.
The hypothesis on the scalar in DJW13 holds since
$B_0(y)=n(\log_p n+1)\ge0$.

In both cases the case $y=0$ deserves its exact value rather than
the rough logarithmic bound.  The polynomial $P_0$ is one, each
$\gamma_{\rho,0}$ and $f_{\rho,0}$ is one, and every gauge
has zero constant term.  Formula \eqref{djw:actualvalue} is
therefore $\mathcal V_0=1$, of valuation zero.  The same argument
shows explicitly that order zero belongs to the admissible set,
as required by \texttt{Windows}.  All systems and coordinates
were fixed before the order was selected.  This proves the exact
window and its original quantifier order.
\end{proof}
% END DJW21

\subsection{The next join and the retained normal-form construction}

The original \texttt{window\_letterDet} theorem first obtains the
normal form from \texttt{RowData.LemmaM}.  The theorem proved here
begins with that exact normal form and constructs every remaining
object explicitly.  Thus, once that original existence theorem is
supplied, DJW21 applies to the same selected letters and proves its
window.  In the alternating setting, an odd shift changes which
side is big.  Hence both selected type spaces have rank at least
$n$, and the first $n$ letters of either original flag give the
injective selections used above.  Their indicator rows are saturated:
the transpose indicators are an integral right inverse.  Renaming the
leading minors as these selected-letter determinants gives the two
original flag versions of \texttt{BiSupply}.

This last description is a conditional assembly, not a claim that
the normal form has now been constructed.  Its original producer
\texttt{RowDataMainTermNormalFormAndLevelOne} uses the adapted
Cartier-chain construction, residue rank estimates, a pairing-kernel
estimate, the weighted gauge remainder, and first-axis vanishing.
Those are the exact upstream tasks retained here.  In particular the
arbitrary series used in the first part has been replaced in the
application by the prescribed restricted gauge, but the common global
theta family and its geometric realization have not been chosen by
this chapter.  The previously closed
\texttt{GLOBAL\_LOCAL\_GAIN} is not reopened or counted as a new
whole closure.

\paragraph{Exact source map.}
All paths below are relative to \path{src/} at the stated source
pin.  The metadata records the exact line ranges and byte hashes.
\begin{itemize}
\item The formal series, integer exponent and derivative conventions
are in
\path{Formalism/IncomingBounds/StatementThmP.lean}, lines 8--42.
The targeted theorem is the distinct, precise
\nolinkurl{Formalism.L2.FormalSeriesBounds.ThmP} in
\path{Formalism/LocalToGlobal/CharacterGerms/}
\path{BiadditiveInSingleVariables.lean}, lines 108--114.
Its row interfaces are at lines 179--224.
\item The window calculation and literal content/window assembly are
in \path{Development/AbelianVarietyTheory/Proofs/UnramifiedDepth/}
\path{SumEqOfIsHomogeneous.lean}, especially lines 771--954 and
961--1085.  The row identity is proved in
\path{BlockValueOnRowsEq.lean}, lines 477--596, using
\path{DetPowDiagonalRowsBlockJetRestrict.lean}.
\item The exact normal-form and restricted-series definitions are in
\path{Formalism/LocalToGlobal/CharacterGerms/}
\path{LatticeInclusionFormsInjective.lean}, lines 83--166.
\path{MainTermImpliesRowHypothesis.lean}, lines 145--194, proves
the full conversion reconstructed in DJW15.
\item The integral completion and differential substitution are in
\path{ConstantCoeffSubstSame.lean}, especially lines 275--350.
The complete actual row construction is
\path{BlockOfNormalForm.lean}, lines 317--434.
The consumer is \path{OddLevelBlockBoundAttained.lean},
lines 89--147, followed by its two flag installations.
\end{itemize}

\paragraph{Published identity and attribution.}
The accepted classical input is S. Caracciolo, A. D. Sokal and
A. Sportiello, \emph{Algebraic/combinatorial proofs of Cayley-type
identities for derivatives of determinants and pfaffians},
Advances in Applied Mathematics \textbf{50} (2013), 474--594,
Theorem 2.1; arXiv:1105.6270, version 2, page 7.
The adjoint-functor reader supplied the useful short organization of
the local window calculation.  The original mathematical argument,
definitions and target results are attributed to Alp\"oge and Claude.
